%% kapa: exam=1 round=35 year=2024 period=1 form=A subject=civil_law subject_ko=민법 seq=20 number=20
%% source: Q-Net 감정평가사 기출문제 (cst00309) artlSeq=5213229 file=raw/1008/5213229/1교시 문제지 원본.pdf
%% conversion: kordoc 4.24.51 -> tools/gen1.py (auto); review=proofread-sonnet; flags=none
\begin{question}{20}
통정허위표시의 무효를 이유로 대항할 수 없는 ‘제3자’에 해당하지 않는 자는? (다툼이 있으면 판례에 따름)
\begin{choices}
  \item 가장소비대차의 계약상의 지위를 이전 받은 자
  \item 가장매매의 목적물에 대하여 저당권을 취득한 자
  \item 가장의 금전소비대차에 기한 대여금채권을 가압류한 자
  \item 가장매매에 의한 매수인으로부터 목적 부동산을 매수하여 소유권이전등기를 마친 자
  \item 가장의 전세권설정계약에 기하여 등기가 마쳐진 전세권에 관하여 저당권을 취득한 자
\end{choices}
\end{question}
\answer{1}
